\documentclass{article}

% ---------------------------------------------------------------------------
% Idioma y codificación
% ---------------------------------------------------------------------------
\usepackage[spanish]{babel}
\usepackage[T1]{fontenc}
\usepackage[utf8]{inputenc}

% ---------------------------------------------------------------------------
% Márgenes y encabezado (formato guía práctica UTA)
% ---------------------------------------------------------------------------
\usepackage[letterpaper,top=2cm,bottom=5cm,left=3cm,right=3cm,marginparwidth=1.75cm]{geometry}
\usepackage{fancyhdr}
\pagestyle{fancy}
\fancyhf{}
\setlength{\headheight}{70pt}

% ---------------------------------------------------------------------------
% Figuras
% ---------------------------------------------------------------------------
\usepackage{graphicx}
\graphicspath{{imagenes/}}
\usepackage{float}

% ---------------------------------------------------------------------------
% Matemáticas y tablas
% ---------------------------------------------------------------------------
\usepackage{amsmath}
\usepackage{amssymb}
\usepackage{booktabs}
\usepackage{array}

% ---------------------------------------------------------------------------
% Enlaces y bibliografía
% ---------------------------------------------------------------------------
\usepackage[colorlinks=true, allcolors=blue]{hyperref}

% ---------------------------------------------------------------------------
% Encabezado institucional (logos + universidad)
% ---------------------------------------------------------------------------
\fancyhead[C]{%
    \begin{minipage}[c]{1.5cm}
        \includegraphics[width=1.5cm]{uta.png}
    \end{minipage}%
    \hfill
    \begin{minipage}[c]{10cm}
        \centering
        \small {\textbf{UNIVERSIDAD TÉCNICA DE AMBATO}} \\[0.1cm]
        \scriptsize \textit{FACULTAD DE INGENIERÍA EN SISTEMAS ELECTRÓNICA E INDUSTRIAL} \\[0.1cm]
        \scriptsize \textbf{CARRERA SOFTWARE} \\[0.1cm]
        \scriptsize \textbf{CICLO ACADÉMICO: JULIO - DICIEMBRE 2026}
        \vspace{0.2cm}
    \end{minipage}%
    \hfill
    \begin{minipage}[c]{1.5cm}
        \centering
        \includegraphics[width=1.5cm]{fisei.png}
    \end{minipage}%
}

\begin{document}

\begin{center}
    {\large \textbf{INFORME DE GUÍA PRÁCTICA}} \\
\end{center}

% ===========================================================================
% PORTADA
% ===========================================================================
\begin{flushleft}
\begin{tabular}{ l p{11cm} }
    \textbf{Tema:} & Optimización del flujo de costo mínimo en la red vial de Huachi (SmartFlow) \\
    \textbf{Unidad de Organización Curricular:} & Profesional \\
    \textbf{Nivel y Paralelo:} & 5 - A \\
    \textbf{Alumnos participante(s):} & Kevin Andres Velasco Huayamave \\
                                      & Garcia Amores Manolo Jose \\
                                      & Juan Pashma \\
    \textbf{Asignatura:} & Investigación Operativa \\
    \textbf{Docente:} & Ing. Ortiz Fernandez William Wladimir \\
\end{tabular}
\end{flushleft}

\tableofcontents
\newpage

% ===========================================================================
% 1. OBJETIVOS
% ===========================================================================
\section{Objetivos}

\subsection{General}
Aplicar la Investigación Operativa a un problema real de transporte, mediante
el desarrollo de \textbf{SmartFlow}: un sistema que modela la red vial del sector
Huachi (Ambato) como un problema de flujo de costo mínimo y determina la
asignación de flujo que minimiza el costo total, respetando la capacidad de las
vías y el balance de oferta y demanda de los nodos.

\subsection{Específicos}
\begin{itemize}
    \item Importar y validar la red vial de Huachi a partir de los datos
          geográficos (GeoJSON) proporcionados por el equipo GIS.
    \item Formular matemáticamente el problema como un modelo de flujo de costo
          mínimo, definiendo conjuntos, parámetros, variable de decisión,
          función objetivo y restricciones.
    \item Implementar la resolución del modelo con el solucionador OR-Tools y
          exponerla mediante una API REST.
    \item Construir escenarios experimentales (base, cierre de vía y caso
          límite) y evaluar el comportamiento del sistema en cada uno.
    \item Interpretar los resultados de la optimización y documentar las
          limitaciones y supuestos del modelo.
\end{itemize}

% ===========================================================================
% 2. MODALIDAD Y DURACIÓN
% ===========================================================================
\section{Modalidad y tiempo de duración}

\subsection{Modalidad}
Presencial.

\subsection{Tiempo de duración}
Presenciales: 4. \\
No presenciales: 0.

% ===========================================================================
% 3. CONTEXTO Y PROBLEMA
% ===========================================================================
\section{Contexto y planteamiento del problema}

El sector de Huachi (Ambato) cuenta con una red de vías que conecta un conjunto
de intersecciones. Cada vía tiene una capacidad limitada (cantidad de flujo que
puede transportar) y un costo unitario asociado a su uso. En la red existen
nodos que \textbf{ofrecen} flujo (oferta) y nodos que \textbf{demandan} flujo
(demanda).

El problema consiste en decidir cuánto flujo enviar por cada vía de manera que:
\begin{enumerate}
    \item se satisfaga el balance de cada nodo (la oferta y la demanda);
    \item no se exceda la capacidad de ninguna vía;
    \item se minimice el costo total de transportar el flujo por la red.
\end{enumerate}

La red se representa como un grafo dirigido, de modo que cada vía puede
recorrerse en uno o ambos sentidos según la información disponible. Cuando una
vía es bidireccional, se modela mediante dos arcos dirigidos, uno por sentido.

% ===========================================================================
% 4. FUNDAMENTO DE INVESTIGACIÓN OPERATIVA
% ===========================================================================
\section{Fundamento de Investigación Operativa}

El modelo seleccionado es el \textbf{problema de flujo de costo mínimo}
(Minimum-Cost Flow Problem, MCFP), uno de los problemas clásicos de la
optimización en redes. La elección se justifica porque el problema cumple,
de forma natural, las tres condiciones que definen un modelo de flujo en red:

\begin{itemize}
    \item \textbf{Estructura de red.} Existe un conjunto de nodos (intersecciones)
          conectados por arcos dirigidos (vías), que forman un grafo.
    \item \textbf{Conservación de flujo.} En cada nodo, el flujo que entra menos
          el flujo que sale debe igualar la oferta (positiva) o la demanda
          (negativa) del nodo; en nodos de tránsito el balance es nulo.
    \item \textbf{Capacidades y costos.} Cada arco tiene una capacidad máxima y
          un costo unitario por enviar una unidad de flujo a través de él.
\end{itemize}

Frente a otras alternativas, el MCFP se ajusta mejor que el problema de la ruta
más corta: este último determina una única ruta de mínimo costo entre dos
puntos, mientras que SmartFlow debe asignar flujo entre múltiples orígenes y
destinos, repartiéndolo por varias vías y respetando capacidades. La
formulación como problema de transporte tampoco basta, pues no representa
directamente las vías intermedias ni sus restricciones de capacidad.

Por lo tanto, el problema se resuelve como un \textbf{problema de flujo de costo
mínimo}, y se optimiza el costo total definido en el modelo (no el tiempo de
viaje ni otros criterios, salvo que estos queden representados en los costos).

% ===========================================================================
% 5. MODELO MATEMÁTICO
% ===========================================================================
\section{Modelo matemático}

\subsection{Conjuntos y parámetros}

Sea $G = (N, A)$ el grafo dirigido de la red, donde:
\begin{itemize}
    \item $N$: conjunto de nodos (intersecciones) de la red.
    \item $A$: conjunto de arcos dirigidos (vías en un sentido determinado).
\end{itemize}

Para cada arco $(i,j) \in A$ se definen:
\begin{itemize}
    \item $c_{ij}$: costo unitario por enviar una unidad de flujo por el arco.
    \item $u_{ij}$: capacidad máxima del arco (límite de flujo).
\end{itemize}

Para cada nodo $i \in N$ se define:
\begin{itemize}
    \item $b_i$: balance del nodo; $b_i > 0$ si es un nodo de oferta,
          $b_i < 0$ si es un nodo de demanda y $b_i = 0$ si es de tránsito.
\end{itemize}

Los balances deben cumplir la condición de factibilidad global:
\begin{equation}
    \sum_{i \in N} b_i = 0 .
\end{equation}

\subsection{Variable de decisión}

La variable de decisión es el flujo enviado por cada arco:
\begin{equation}
    x_{ij} \ge 0 \quad \forall\, (i,j) \in A .
\end{equation}

\subsection{Función objetivo}

Se minimiza el costo total de transportar el flujo por la red:
\begin{equation}
    \min \; \sum_{(i,j) \in A} c_{ij} \, x_{ij} .
\end{equation}

\subsection{Restricciones}

\textbf{Conservación de flujo.} Para cada nodo $i \in N$, la diferencia entre
el flujo que sale y el que entra debe igualar su balance:
\begin{equation}
    \sum_{j : (i,j) \in A} x_{ij} \; - \sum_{j : (j,i) \in A} x_{ji} \;=\; b_i
    \quad \forall\, i \in N .
\end{equation}

\textbf{Capacidades.} El flujo por cada arco no puede exceder su capacidad:
\begin{equation}
    0 \le x_{ij} \le u_{ij} \quad \forall\, (i,j) \in A .
\end{equation}

El problema es un programa lineal (función objetivo y restricciones lineales).
Su solución, cuando existe, indica cuánto flujo debe circular por cada vía para
minimizar el costo total respetando oferta, demanda y capacidades. Si no existe
ninguna asignación que satisfaga todas las restricciones, el problema es
\textbf{infactible}.

% ===========================================================================
% 6. DATOS, SUPUESTOS Y ALCANCE
% ===========================================================================
\section{Datos, supuestos y alcance}

\subsection{Datos geográficos}
La topología de la red (nodos y vías) proviene de los datos de OpenStreetMap
procesados por el equipo GIS. Se importan como GeoJSON: los nodos son puntos y
las vías son líneas. El importador valida que la estructura sea consistente
(nodos con identificador único, vías cuyos extremos coinciden con nodos
existentes, etc.).

\subsection{Supuestos experimentales}
Los valores de capacidad, costo unitario y oferta/demanda de los escenarios son
\textbf{supuestos experimentales} definidos por el equipo con fines
académicos, no mediciones de campo. El costo unitario representa el costo
definido para el modelo, y su interpretación (por ejemplo, longitud o un costo
asociado) debe declararse de forma explícita al presentar los resultados.

\subsection{Alcance}
SmartFlow optimiza la asignación de flujo para el costo definido. No debe
interpretarse el resultado como la ruta más rápida ni más segura salvo que
dichos criterios estén efectivamente representados en los costos del modelo.

% ===========================================================================
% 7. METODOLOGÍA Y FUNCIONAMIENTO DEL SISTEMA
% ===========================================================================
\section{Metodología y funcionamiento del sistema}

\subsection{Arquitectura}
El sistema se compone de:
\begin{itemize}
    \item \textbf{Frontend}: visualización de la red en el mapa y selección de
          escenarios.
    \item \textbf{Backend (API REST)}: validación de datos, construcción de
          escenarios y ejecución del optimizador.
    \item \textbf{Motor de optimización}: OR-Tools, que resuelve el problema de
          flujo de costo mínimo.
    \item \textbf{Datos geográficos}: GeoJSON procesado por el equipo GIS.
\end{itemize}

\subsection{Proceso de optimización}
\begin{enumerate}
    \item Se importa y valida la red geográfica (nodos y vías).
    \item Se carga la semilla del escenario (balances, capacidades, costos y
          configuraciones como cierres de vía).
    \item Se construye el modelo de flujo de costo mínimo a partir de la red y
          el escenario.
    \item Se resuelve con OR-Tools.
    \item Se devuelve el estado (óptimo o infactible), el costo total y el
          flujo por cada vía.
\end{enumerate}

% ===========================================================================
% 8. MANUAL DE USUARIO
% ===========================================================================
\section{Manual de usuario}

\subsection{Inicio del sistema}
\begin{enumerate}
    \item Iniciar el backend (API) y el frontend según las instrucciones del
          proyecto.
    \item Abrir la aplicación en el navegador para visualizar el mapa de la red
          de Huachi.
\end{enumerate}

\subsection{Selección y ejecución de un escenario}
\begin{enumerate}
    \item Seleccionar uno de los escenarios disponibles: \textbf{Base},
          \textbf{Cierre de vía} o \textbf{Caso límite}.
    \item Ejecutar la optimización.
\end{enumerate}

\subsection{Interpretación de resultados}
\begin{itemize}
    \item \textbf{Estado}: indica si existe solución (\texttt{optimal}) o si el
          problema es \texttt{infeasible}.
    \item \textbf{Costo total}: valor de la función objetivo en el óptimo.
    \item \textbf{Flujo por vía}: cantidad de flujo asignada a cada vía.
    \item \textbf{Capacidad y utilización}: relación entre el flujo asignado y
          la capacidad de la vía.
\end{itemize}

Un resultado \texttt{infeasible} significa que, con las restricciones dadas
(por ejemplo, capacidades reducidas o vías cerradas), no es posible satisfacer
la oferta y la demanda de la red.

% ===========================================================================
% 9. ACTIVIDADES Y ESCENARIOS EVALUADOS
% ===========================================================================
\section{Actividades por desarrollar}

Se desarrolla la guía práctica mediante la implementación del sistema SmartFlow
y la evaluación de tres escenarios:

\begin{itemize}
    \item \textbf{Base}: todas las vías habilitadas con sus capacidades y
          costos originales.
    \item \textbf{Cierre de vía}: se deshabilitan una o más vías (simulación de
          un cierre) y se observa cómo cambia la asignación de flujo.
    \item \textbf{Caso límite}: se reducen las capacidades mediante un factor,
          llevando el sistema a una situación cercana o directamente a la
          infactibilidad.
\end{itemize}

% ===========================================================================
% 10. RESULTADOS OBTENIDOS
% ===========================================================================
\section{Resultados obtenidos}

\subsection{Resultados de la optimización}
La siguiente tabla resume los resultados observados para la red de Huachi
(pendiente de contrastar con la versión final antes de la entrega):

\begin{center}
\begin{tabular}{ l c c }
    \toprule
    \textbf{Escenario} & \textbf{Estado} & \textbf{Costo total} \\
    \midrule
    Base            & optimal    & 34\,080 \\
    Cierre de vía   & optimal    & 42\,600 \\
    Caso límite     & infeasible & --- \\
    \bottomrule
\end{tabular}
\end{center}

\subsection{Validación}
La implementación se valida mediante pruebas automatizadas que verifican el
importador de datos, la construcción de escenarios, la resolución del modelo y
la persistencia de resultados. Los valores numéricos se restringen al rango de
enteros seguros de JavaScript para garantizar consistencia entre el frontend y
el backend.

% ===========================================================================
% 11. HABILIDADES BLANDAS
% ===========================================================================
\section{Habilidades blandas empleadas en la práctica}

\begin{itemize}
    \item[$\square$] Liderazgo
    \item[$\blacksquare$] Trabajo en equipo
    \item[$\square$] Comunicación asertiva
    \item[$\square$] La empatía
    \item[$\blacksquare$] Pensamiento crítico
    \item[$\square$] Flexibilidad
    \item[$\square$] La resolución de conflictos
    \item[$\square$] Adaptabilidad
    \item[$\blacksquare$] Responsabilidad
\end{itemize}

% ===========================================================================
% 12. CONCLUSIONES
% ===========================================================================
\section{Conclusiones}

\begin{itemize}
    \item El problema de flujo de costo mínimo permite modelar de forma
          adecuada la asignación de flujo sobre la red vial de Huachi, al
          representar nodos, vías, capacidades, oferta y demanda.
    \item La separación entre datos geográficos y supuestos experimentales es
          clave: los resultados son válidos para el modelo definido y deben
          interpretarse en función de cómo se definan los costos.
    \item La implementación con OR-Tools resuelve el modelo de forma eficiente
          y permite evaluar escenarios distintos (base, cierre y caso límite).
    \item Los casos en que no existe solución (infactibles) son tan
          informativos como los óptimos, pues evidencian restricciones
          incompatibles en la red.
\end{itemize}

% ===========================================================================
% 13. RECOMENDACIONES
% ===========================================================================
\section{Recomendaciones}

\begin{itemize}
    \item Declarar de forma explícita el significado de los costos y de los
          supuestos experimentales al presentar los resultados.
    \item Validar la geometría y la topología de la red con el equipo GIS
          (sentidos de circulación, vías paralelas y bucles).
    \item Ampliar los escenarios para estudiar el impacto de distintos cierres
          y variaciones de capacidad sobre la red.
    \item Documentar con capturas cada escenario para facilitar la
          interpretación de los resultados por parte del lector.
\end{itemize}

% ===========================================================================
% 14. REFERENCIAS
% ===========================================================================
\section{Referencias}

\nocite{*}
\bibliographystyle{ieeetr}
\bibliography{referencias}

\end{document}
